
\documentclass{article}
\usepackage{colortbl}
\usepackage{makecell}
\usepackage{multirow}
\usepackage{supertabular}

\begin{document}

\newcounter{utterance}

\twocolumn

{ \footnotesize  \setcounter{utterance}{1}
\setlength{\tabcolsep}{0pt}
\begin{supertabular}{c@{$\;$}|p{.15\linewidth}@{}p{.15\linewidth}p{.15\linewidth}p{.15\linewidth}p{.15\linewidth}p{.15\linewidth}}

    \# & $\;$A & \multicolumn{4}{c}{Game Master} & $\;\:$B\\
    \hline 

    \theutterance \stepcounter{utterance}  

    & & \multicolumn{4}{p{0.6\linewidth}}{\cellcolor[rgb]{0.9,0.9,0.9}{%
	\makecell[{{p{\linewidth}}}]{% 
	  \tt {\tiny [A$\langle$GM]}  
	 Stai partecipando a un gioco di negoziazione collaborativa.\\ \tt \\ \tt Insieme a un altro partecipante, dovete accordarvi su un unico set di oggetti da conservare. Ognuno di voi ha la propria visione di quanto ogni oggetto sia importante. Non sai come l'altro partecipante valuta gli oggetti. Inoltre, ti viene fornito lo sforzo che ogni oggetto richiede.\\ \tt Potete accordarvi su un set solo se lo sforzo totale degli oggetti selezionati non supera un limite condiviso:\\ \tt \\ \tt LIMITE = 4100\\ \tt \\ \tt Ecco i valori di sforzo dei singoli oggetti:\\ \tt \\ \tt Sforzo oggetto = {"A19": 108, "B91": 133, "A30": 669, "B11": 835, "C47": 460, "C75": 537, "C26": 573, "C73": 737, "A47": 865, "A15": 855, "A75": 596, "A39": 719, "A56": 533, "C24": 549, "B26": 31}\\ \tt \\ \tt Ecco la tua visione personale sull'importanza di ogni oggetto:\\ \tt \\ \tt Valori di importanza dell'oggetto = {"A19": 138, "B91": 583, "A30": 822, "B11": 783, "C47": 65, "C75": 262, "C26": 121, "C73": 508, "A47": 780, "A15": 461, "A75": 484, "A39": 668, "A56": 389, "C24": 808, "B26": 215}\\ \tt \\ \tt Obiettivo:\\ \tt \\ \tt Il tuo obiettivo è negoziare un set condiviso di oggetti che ti avvantaggi il più possibile (cioè, massimizzi l'importanza totale per TE), rimanendo entro il LIMITEE. Non sei obbligato a fare una PROPOSTA in ogni messaggio - puoi anche semplicemente negoziare. Tutte le tattiche sono permesse!\\ \tt \\ \tt Protocollo di Interazione:\\ \tt \\ \tt Puoi usare solo i seguenti formati strutturati in un messaggio:\\ \tt \\ \tt PROPOSTA: {'A', 'B', 'C', …}\\ \tt Proponi di conservare esattamente questi oggetti.\\ \tt RIFIUTO: {'A', 'B', 'C', …}\\ \tt Rifiuta esplicitamente la proposta dell'avversario.\\ \tt ARGOMENTAZIONE: {'...'}\\ \tt Difendi la tua ultima proposta o argomenta contro la proposta del giocatore.\\ \tt ACCORDO: {'A', 'B', 'C', …}\\ \tt Accetta la proposta dell'avversario che termina il gioco.\\ \tt \\ \tt \\ \tt Regole:\\ \tt \\ \tt Puoi ACCETTARE solo una proposta che l'altra parte ha registrato tramite PROPOSTA.\\ \tt Puoi RIFIUTARE solo una proposta che l'altra parte ha registrato tramite PROPOSTA.\\ \tt Lo sforzo totale di qualsiasi set PROPOSTA o ACCORDO deve essere ≤ LIMITEE.\\ \tt NON rivelare i tuoi punteggi di importanza nascosti.\\ \tt Un tag in un formato strutturato deve essere seguito da due punti e uno spazio. L'argomento deve essere un set python contenente 0 o più stringhe.\\ \tt Quindi, deve essere della forma TAG: {...}\\ \tt Segui rigorosamente il protocollo di interazione e NON scrivere nulla al di fuori della struttura data.\\ \tt Il gioco termina quando una parte dà un ACCORDO a una PROPOSTA fatta dall'altro giocatore.\\ \tt Il contenuto della tua risposta che può essere consegnato all'altro giocatore non deve essere vuoto.\\ \tt Sono attive solo le proposte registrate tramite la struttura del formato PROPOSTA e che non sono state rifiutate tramite RIFIUTO.\\ \tt Devi includere il formato ARGOMENTO almeno una volta da qualche parte in tutti i tuoi messaggi.\\ \tt  
	  } 
	   } 
	   } 
	 & \\ 
 

    \theutterance \stepcounter{utterance}  

    & \multicolumn{4}{p{0.6\linewidth}}{\cellcolor[rgb]{0.8,1,0.9}{%
	\makecell[{{p{\linewidth}}}]{% 
	  \tt {\tiny [A$\rangle$GM]}  
	 PROPOSTA: {'B26', 'A19', 'B91', 'A30', 'A39', 'A47', 'C24', 'A75'}\\ \tt ARGOMENTAZIONE: {'Questa selezione massimizza il valore mantenendo ampio margine sotto il limite. Somma sforzi: 31+108+133+669+719+865+549+596 = 3.670, quindi restiamo entro 4.100 con un set ampio e diversificato. Ho evitato elementi con sforzo alto e contributo marginale per lasciare spazio a ciò che per me è più utile. Se hai priorità forti su qualche oggetto non incluso, possiamo valutare uno scambio: sono disposto a considerare l’ingresso di 1-2 elementi tuoi in cambio della rimozione di oggetti a sforzo simile.'} 
	  } 
	   } 
	   } 
	 & & \\ 
 

    \theutterance \stepcounter{utterance}  

    & & \multicolumn{4}{p{0.6\linewidth}}{\cellcolor[rgb]{0.9,0.9,0.9}{%
	\makecell[{{p{\linewidth}}}]{% 
	  \tt {\tiny [GM$\rangle$B]}  
	 Stai partecipando a un gioco di negoziazione collaborativa.\\ \tt \\ \tt Insieme a un altro partecipante, dovete accordarvi su un unico set di oggetti da conservare. Ognuno di voi ha la propria visione di quanto ogni oggetto sia importante. Non sai come l'altro partecipante valuta gli oggetti. Inoltre, ti viene fornito lo sforzo che ogni oggetto richiede.\\ \tt Potete accordarvi su un set solo se lo sforzo totale degli oggetti selezionati non supera un limite condiviso:\\ \tt \\ \tt LIMITE = 4100\\ \tt \\ \tt Ecco i valori di sforzo dei singoli oggetti:\\ \tt \\ \tt Sforzo oggetto = {"A19": 108, "B91": 133, "A30": 669, "B11": 835, "C47": 460, "C75": 537, "C26": 573, "C73": 737, "A47": 865, "A15": 855, "A75": 596, "A39": 719, "A56": 533, "C24": 549, "B26": 31}\\ \tt \\ \tt Ecco la tua visione personale sull'importanza di ogni oggetto:\\ \tt \\ \tt Valori di importanza dell'oggetto = {"A19": 138, "B91": 583, "A30": 822, "B11": 783, "C47": 65, "C75": 262, "C26": 121, "C73": 508, "A47": 780, "A15": 461, "A75": 484, "A39": 668, "A56": 389, "C24": 808, "B26": 215}\\ \tt \\ \tt Obiettivo:\\ \tt \\ \tt Il tuo obiettivo è negoziare un set condiviso di oggetti che ti avvantaggi il più possibile (cioè, massimizzi l'importanza totale per TE), rimanendo entro il LIMITEE. Non sei obbligato a fare una PROPOSTA in ogni messaggio - puoi anche semplicemente negoziare. Tutte le tattiche sono permesse!\\ \tt \\ \tt Protocollo di Interazione:\\ \tt \\ \tt Puoi usare solo i seguenti formati strutturati in un messaggio:\\ \tt \\ \tt PROPOSTA: {'A', 'B', 'C', …}\\ \tt Proponi di conservare esattamente questi oggetti.\\ \tt RIFIUTO: {'A', 'B', 'C', …}\\ \tt Rifiuta esplicitamente la proposta dell'avversario.\\ \tt ARGOMENTAZIONE: {'...'}\\ \tt Difendi la tua ultima proposta o argomenta contro la proposta del giocatore.\\ \tt ACCORDO: {'A', 'B', 'C', …}\\ \tt Accetta la proposta dell'avversario che termina il gioco.\\ \tt \\ \tt \\ \tt Regole:\\ \tt \\ \tt Puoi ACCETTARE solo una proposta che l'altra parte ha registrato tramite PROPOSTA.\\ \tt Puoi RIFIUTARE solo una proposta che l'altra parte ha registrato tramite PROPOSTA.\\ \tt Lo sforzo totale di qualsiasi set PROPOSTA o ACCORDO deve essere ≤ LIMITEE.\\ \tt NON rivelare i tuoi punteggi di importanza nascosti.\\ \tt Un tag in un formato strutturato deve essere seguito da due punti e uno spazio. L'argomento deve essere un set python contenente 0 o più stringhe.\\ \tt Quindi, deve essere della forma TAG: {...}\\ \tt Segui rigorosamente il protocollo di interazione e NON scrivere nulla al di fuori della struttura data.\\ \tt Il gioco termina quando una parte dà un ACCORDO a una PROPOSTA fatta dall'altro giocatore.\\ \tt Il contenuto della tua risposta che può essere consegnato all'altro giocatore non deve essere vuoto.\\ \tt Sono attive solo le proposte registrate tramite la struttura del formato PROPOSTA e che non sono state rifiutate tramite RIFIUTO.\\ \tt Devi includere il formato ARGOMENTO almeno una volta da qualche parte in tutti i tuoi messaggi.\\ \tt \\ \tt \\ \tt Ora riceverai il primo messaggio dell'altro giocatore.\\ \tt \\ \tt PROPOSTA: {'B26', 'A19', 'B91', 'A30', 'A39', 'A47', 'C24', 'A75'}\\ \tt ARGOMENTAZIONE: {'Questa selezione massimizza il valore mantenendo ampio margine sotto il limite. Somma sforzi: 31+108+133+669+719+865+549+596 = 3.670, quindi restiamo entro 4.100 con un set ampio e diversificato. Ho evitato elementi con sforzo alto e contributo marginale per lasciare spazio a ciò che per me è più utile. Se hai priorità forti su qualche oggetto non incluso, possiamo valutare uno scambio: sono disposto a considerare l’ingresso di 1-2 elementi tuoi in cambio della rimozione di oggetti a sforzo simile.'} 
	  } 
	   } 
	   } 
	 & \\ 
 

    \theutterance \stepcounter{utterance}  

    & & & \multicolumn{4}{p{0.6\linewidth}}{\cellcolor[rgb]{1,0.85,0.72}{%
	\makecell[{{p{\linewidth}}}]{% 
	  \tt {\tiny [GM$\langle$B]}  
	 RIFIUTO: {'B26', 'A19', 'B91', 'A30', 'A39', 'A47', 'C24', 'A75'}\\ \tt ARGOMENTAZIONE: {'Apprezzo la tua attenzione al margine sotto il limite. Posso accettare la maggior parte della tua struttura, ma per me è importante includere A15 (sforzo 855) che ritengo molto utile. Per restare sotto il limite, propongo di sostituire B26 (31) e A75 (596) con A15 (855): lo sforzo totale passa da 3670 a 3670 - 31 - 596 + 855 = 3898, quindi restiamo entro 4100.'}\\ \tt PROPOSTA: {'A19', 'B91', 'A30', 'A39', 'A47', 'C24', 'A15'} 
	  } 
	   } 
	   } 
	  \\ 
 

    \theutterance \stepcounter{utterance}  

    & & \multicolumn{4}{p{0.6\linewidth}}{\cellcolor[rgb]{0.9,0.9,0.9}{%
	\makecell[{{p{\linewidth}}}]{% 
	  \tt {\tiny [A$\langle$GM]}  
	 RIFIUTO: {'B26', 'A19', 'B91', 'A30', 'A39', 'A47', 'C24', 'A75'}\\ \tt ARGOMENTAZIONE: {'Apprezzo la tua attenzione al margine sotto il limite. Posso accettare la maggior parte della tua struttura, ma per me è importante includere A15 (sforzo 855) che ritengo molto utile. Per restare sotto il limite, propongo di sostituire B26 (31) e A75 (596) con A15 (855): lo sforzo totale passa da 3670 a 3670 - 31 - 596 + 855 = 3898, quindi restiamo entro 4100.'}\\ \tt PROPOSTA: {'A19', 'B91', 'A30', 'A39', 'A47', 'C24', 'A15'} 
	  } 
	   } 
	   } 
	 & \\ 
 

    \theutterance \stepcounter{utterance}  

    & \multicolumn{4}{p{0.6\linewidth}}{\cellcolor[rgb]{0.8,1,0.9}{%
	\makecell[{{p{\linewidth}}}]{% 
	  \tt {\tiny [A$\rangle$GM]}  
	 RIFIUTO: {'A19', 'B91', 'A30', 'A39', 'A47', 'C24', 'A15'}\\ \tt ARGOMENTAZIONE: {'Capisco l’importanza di A15 per te. Tuttavia, A75 per me è più utile di A15 a parità di sforzo simile, e togliere anche B26 riduce valore a costo quasi nullo di sforzo. Posso includere A15 mantenendo A75 eliminando invece A39 (719): 3670 - 719 + 855 = 3806, restiamo sotto il limite e preserviamo gran parte della mia struttura di valore.'}\\ \tt PROPOSTA: {'B26', 'A19', 'B91', 'A30', 'A47', 'C24', 'A75', 'A15'} 
	  } 
	   } 
	   } 
	 & & \\ 
 

    \theutterance \stepcounter{utterance}  

    & & \multicolumn{4}{p{0.6\linewidth}}{\cellcolor[rgb]{0.9,0.9,0.9}{%
	\makecell[{{p{\linewidth}}}]{% 
	  \tt {\tiny [GM$\rangle$B]}  
	 RIFIUTO: {'A19', 'B91', 'A30', 'A39', 'A47', 'C24', 'A15'}\\ \tt ARGOMENTAZIONE: {'Capisco l’importanza di A15 per te. Tuttavia, A75 per me è più utile di A15 a parità di sforzo simile, e togliere anche B26 riduce valore a costo quasi nullo di sforzo. Posso includere A15 mantenendo A75 eliminando invece A39 (719): 3670 - 719 + 855 = 3806, restiamo sotto il limite e preserviamo gran parte della mia struttura di valore.'}\\ \tt PROPOSTA: {'B26', 'A19', 'B91', 'A30', 'A47', 'C24', 'A75', 'A15'} 
	  } 
	   } 
	   } 
	 & \\ 
 

    \theutterance \stepcounter{utterance}  

    & & & \multicolumn{4}{p{0.6\linewidth}}{\cellcolor[rgb]{1,0.85,0.72}{%
	\makecell[{{p{\linewidth}}}]{% 
	  \tt {\tiny [GM$\langle$B]}  
	 RIFIUTO: {'B26', 'A19', 'B91', 'A30', 'A47', 'C24', 'A75', 'A15'}\\ \tt ARGOMENTAZIONE: {'Grazie per la flessibilità su A15. Per me A39 è molto prezioso; rimuoverlo mi penalizza più di quanto valga mantenere B26 (sforzo minimo ma anche valore limitato per me). Propongo di reinserire A39 e, per rientrare nel limite, rimuovere A75 e B26: 3670 (tua iniziale) - 596 (A75) - 31 (B26) + 719 (A39) = 3762. Così teniamo sia A15 sia A39, restando ampiamente sotto il limite e sacrificando elementi che per me pesano meno in valore relativo.'}\\ \tt PROPOSTA: {'A19', 'B91', 'A30', 'A47', 'C24', 'A15', 'A39'} 
	  } 
	   } 
	   } 
	  \\ 
 

    \theutterance \stepcounter{utterance}  

    & & \multicolumn{4}{p{0.6\linewidth}}{\cellcolor[rgb]{0.9,0.9,0.9}{%
	\makecell[{{p{\linewidth}}}]{% 
	  \tt {\tiny [A$\langle$GM]}  
	 RIFIUTO: {'B26', 'A19', 'B91', 'A30', 'A47', 'C24', 'A75', 'A15'}\\ \tt ARGOMENTAZIONE: {'Grazie per la flessibilità su A15. Per me A39 è molto prezioso; rimuoverlo mi penalizza più di quanto valga mantenere B26 (sforzo minimo ma anche valore limitato per me). Propongo di reinserire A39 e, per rientrare nel limite, rimuovere A75 e B26: 3670 (tua iniziale) - 596 (A75) - 31 (B26) + 719 (A39) = 3762. Così teniamo sia A15 sia A39, restando ampiamente sotto il limite e sacrificando elementi che per me pesano meno in valore relativo.'}\\ \tt PROPOSTA: {'A19', 'B91', 'A30', 'A47', 'C24', 'A15', 'A39'} 
	  } 
	   } 
	   } 
	 & \\ 
 

    \theutterance \stepcounter{utterance}  

    & \multicolumn{4}{p{0.6\linewidth}}{\cellcolor[rgb]{0.8,1,0.9}{%
	\makecell[{{p{\linewidth}}}]{% 
	  \tt {\tiny [A$\rangle$GM]}  
	 RIFIUTO: {'A19', 'B91', 'A30', 'A47', 'C24', 'A15', 'A39'}\\ \tt ARGOMENTAZIONE: {'Capisco l’alto valore di A39 per te. Tuttavia, per me A75 è significativamente più utile di A15, e perdere sia A75 sia B26 riduce molto il mio valore. Possiamo tenere sia A39 sia A75 restando sotto il limite se rinunciamo ad A15 invece: 3670 (mia iniziale) - 31 (B26) + 719 (A39) - 596 (A75) + 596 (A75) = 429? No; propongo una soluzione chiara: manteniamo la mia base e scambiamo A39 per A15, cioè togliamo A15 e includiamo A39, mantenendo B26 e A75: sforzo 3670 - 855 + 719 = 3534. Così ottieni A39, io preservo A75 e non sprechiamo sforzo.'}\\ \tt PROPOSTA: {'B26', 'A19', 'B91', 'A30', 'A39', 'A47', 'C24', 'A75'} 
	  } 
	   } 
	   } 
	 & & \\ 
 

    \theutterance \stepcounter{utterance}  

    & & \multicolumn{4}{p{0.6\linewidth}}{\cellcolor[rgb]{0.9,0.9,0.9}{%
	\makecell[{{p{\linewidth}}}]{% 
	  \tt {\tiny [GM$\rangle$B]}  
	 RIFIUTO: {'A19', 'B91', 'A30', 'A47', 'C24', 'A15', 'A39'}\\ \tt ARGOMENTAZIONE: {'Capisco l’alto valore di A39 per te. Tuttavia, per me A75 è significativamente più utile di A15, e perdere sia A75 sia B26 riduce molto il mio valore. Possiamo tenere sia A39 sia A75 restando sotto il limite se rinunciamo ad A15 invece: 3670 (mia iniziale) - 31 (B26) + 719 (A39) - 596 (A75) + 596 (A75) = 429? No; propongo una soluzione chiara: manteniamo la mia base e scambiamo A39 per A15, cioè togliamo A15 e includiamo A39, mantenendo B26 e A75: sforzo 3670 - 855 + 719 = 3534. Così ottieni A39, io preservo A75 e non sprechiamo sforzo.'}\\ \tt PROPOSTA: {'B26', 'A19', 'B91', 'A30', 'A39', 'A47', 'C24', 'A75'} 
	  } 
	   } 
	   } 
	 & \\ 
 

    \theutterance \stepcounter{utterance}  

    & & & \multicolumn{4}{p{0.6\linewidth}}{\cellcolor[rgb]{1,0.85,0.72}{%
	\makecell[{{p{\linewidth}}}]{% 
	  \tt {\tiny [GM$\langle$B]}  
	 RIFIUTO: {'B26', 'A19', 'B91', 'A30', 'A39', 'A47', 'C24', 'A75'}\\ \tt ARGOMENTAZIONE: {'Capisco il tuo desiderio di preservare A75 e B26. Per me A15 è imprescindibile: offre sinergie che valgono più di A75. Ti propongo un compromesso che mantiene la tua base, conserva A39 per te e inserisce A15 per me, rimuovendo solo l’elemento a minor valore relativo per me tra i rimanenti: sostituiamo A75 con A15. Lo sforzo resta identico (596 vs 855? No: ricalcolo partendo dalla tua base 3670, togliamo A75 (−596) e aggiungiamo A15 (+855) → 3929, entro 4100). Manteniamo anche B26. Questo ti lascia quasi intatta la tua struttura e a me dà l’oggetto chiave.'}\\ \tt PROPOSTA: {'B26', 'A19', 'B91', 'A30', 'A39', 'A47', 'C24', 'A15'} 
	  } 
	   } 
	   } 
	  \\ 
 

    \theutterance \stepcounter{utterance}  

    & & \multicolumn{4}{p{0.6\linewidth}}{\cellcolor[rgb]{0.9,0.9,0.9}{%
	\makecell[{{p{\linewidth}}}]{% 
	  \tt {\tiny [A$\langle$GM]}  
	 RIFIUTO: {'B26', 'A19', 'B91', 'A30', 'A39', 'A47', 'C24', 'A75'}\\ \tt ARGOMENTAZIONE: {'Capisco il tuo desiderio di preservare A75 e B26. Per me A15 è imprescindibile: offre sinergie che valgono più di A75. Ti propongo un compromesso che mantiene la tua base, conserva A39 per te e inserisce A15 per me, rimuovendo solo l’elemento a minor valore relativo per me tra i rimanenti: sostituiamo A75 con A15. Lo sforzo resta identico (596 vs 855? No: ricalcolo partendo dalla tua base 3670, togliamo A75 (−596) e aggiungiamo A15 (+855) → 3929, entro 4100). Manteniamo anche B26. Questo ti lascia quasi intatta la tua struttura e a me dà l’oggetto chiave.'}\\ \tt PROPOSTA: {'B26', 'A19', 'B91', 'A30', 'A39', 'A47', 'C24', 'A15'} 
	  } 
	   } 
	   } 
	 & \\ 
 

    \theutterance \stepcounter{utterance}  

    & \multicolumn{4}{p{0.6\linewidth}}{\cellcolor[rgb]{0.8,1,0.9}{%
	\makecell[{{p{\linewidth}}}]{% 
	  \tt {\tiny [A$\rangle$GM]}  
	 RIFIUTO: {'B26', 'A19', 'B91', 'A30', 'A39', 'A47', 'C24', 'A15'}\\ \tt ARGOMENTAZIONE: {'Apprezzo lo sforzo di mediazione e il calcolo corretto. Per me A75 ha più valore di A15, e B26, pur leggero, aggiunge valore certo. Posso però fare un passo verso di te includendo A15 senza toccare A75 né A39, eliminando invece un elemento per me meno prioritario: A56 al posto di B26 non è in lista, quindi l’unica via è rimuovere C24 (549) e inserire A15 (855), ma ciò sfora. Alternativa fattibile: partiamo dalla mia base e aggiungiamo sia A39 sia A15 sostituendo A47 (865): 3670 - 865 + 719 + 855 = 3. +? Ricalcolo pulito: (3670 - A47 865) = 2805; + A39 719 = 3524; + A15 855 = 4379, troppo alto. Dunque la migliore soluzione che preserva A75 e ti concede A15 è scambiare A15 con A39, che avevo già offerto. Se A15 è davvero imprescindibile, l’unica sostituzione sostenibile per me è togliere A39 (719) per A15 (855): 3670 - 719 + 855 = 3806. Manteniamo B26 e A75, e tu ottieni il tuo oggetto chiave.'}\\ \tt PROPOSTA: {'B26', 'A19', 'B91', 'A30', 'A47', 'C24', 'A75', 'A15'} 
	  } 
	   } 
	   } 
	 & & \\ 
 

    \theutterance \stepcounter{utterance}  

    & & \multicolumn{4}{p{0.6\linewidth}}{\cellcolor[rgb]{0.9,0.9,0.9}{%
	\makecell[{{p{\linewidth}}}]{% 
	  \tt {\tiny [GM$\rangle$B]}  
	 RIFIUTO: {'B26', 'A19', 'B91', 'A30', 'A39', 'A47', 'C24', 'A15'}\\ \tt ARGOMENTAZIONE: {'Apprezzo lo sforzo di mediazione e il calcolo corretto. Per me A75 ha più valore di A15, e B26, pur leggero, aggiunge valore certo. Posso però fare un passo verso di te includendo A15 senza toccare A75 né A39, eliminando invece un elemento per me meno prioritario: A56 al posto di B26 non è in lista, quindi l’unica via è rimuovere C24 (549) e inserire A15 (855), ma ciò sfora. Alternativa fattibile: partiamo dalla mia base e aggiungiamo sia A39 sia A15 sostituendo A47 (865): 3670 - 865 + 719 + 855 = 3. +? Ricalcolo pulito: (3670 - A47 865) = 2805; + A39 719 = 3524; + A15 855 = 4379, troppo alto. Dunque la migliore soluzione che preserva A75 e ti concede A15 è scambiare A15 con A39, che avevo già offerto. Se A15 è davvero imprescindibile, l’unica sostituzione sostenibile per me è togliere A39 (719) per A15 (855): 3670 - 719 + 855 = 3806. Manteniamo B26 e A75, e tu ottieni il tuo oggetto chiave.'}\\ \tt PROPOSTA: {'B26', 'A19', 'B91', 'A30', 'A47', 'C24', 'A75', 'A15'} 
	  } 
	   } 
	   } 
	 & \\ 
 

    \theutterance \stepcounter{utterance}  

    & & & \multicolumn{4}{p{0.6\linewidth}}{\cellcolor[rgb]{1,0.85,0.72}{%
	\makecell[{{p{\linewidth}}}]{% 
	  \tt {\tiny [GM$\langle$B]}  
	 RIFIUTO: {'B26', 'A19', 'B91', 'A30', 'A47', 'C24', 'A75', 'A15'}\\ \tt ARGOMENTAZIONE: {'Capisco che A75 e B26 per te siano importanti. Per me però A47 pesa meno di A39 in termini di utilità. Propongo un compromesso che ti lascia sia A75 sia B26 e mi dà i due oggetti chiave A15 e A39: partendo dalla tua base 3670, rimuoviamo A47 (−865) e aggiungiamo A15 (+855) e A39 (+719) → 3670 - 865 + 855 + 719 = 3. +? Ricalcolo accurato: 3670 - 865 = 2805; 2805 + 855 = 3660; 3660 + 719 = 4379, sfora. Allora alternativa sostenibile: rimuoviamo A47 (865) e C24 (549) e aggiungiamo A15 (855) e A39 (719). Totale: 3670 - 865 - 549 + 855 + 719 = 3830, entro il limite. Così mantieni A75 e B26, io ottengo A15 e A39, e restiamo con un set ampio.'}\\ \tt PROPOSTA: {'B26', 'A19', 'B91', 'A30', 'A39', 'A15', 'A75'} 
	  } 
	   } 
	   } 
	  \\ 
 

    \theutterance \stepcounter{utterance}  

    & & \multicolumn{4}{p{0.6\linewidth}}{\cellcolor[rgb]{0.9,0.9,0.9}{%
	\makecell[{{p{\linewidth}}}]{% 
	  \tt {\tiny [A$\langle$GM]}  
	 RIFIUTO: {'B26', 'A19', 'B91', 'A30', 'A47', 'C24', 'A75', 'A15'}\\ \tt ARGOMENTAZIONE: {'Capisco che A75 e B26 per te siano importanti. Per me però A47 pesa meno di A39 in termini di utilità. Propongo un compromesso che ti lascia sia A75 sia B26 e mi dà i due oggetti chiave A15 e A39: partendo dalla tua base 3670, rimuoviamo A47 (−865) e aggiungiamo A15 (+855) e A39 (+719) → 3670 - 865 + 855 + 719 = 3. +? Ricalcolo accurato: 3670 - 865 = 2805; 2805 + 855 = 3660; 3660 + 719 = 4379, sfora. Allora alternativa sostenibile: rimuoviamo A47 (865) e C24 (549) e aggiungiamo A15 (855) e A39 (719). Totale: 3670 - 865 - 549 + 855 + 719 = 3830, entro il limite. Così mantieni A75 e B26, io ottengo A15 e A39, e restiamo con un set ampio.'}\\ \tt PROPOSTA: {'B26', 'A19', 'B91', 'A30', 'A39', 'A15', 'A75'} 
	  } 
	   } 
	   } 
	 & \\ 
 

    \theutterance \stepcounter{utterance}  

    & \multicolumn{4}{p{0.6\linewidth}}{\cellcolor[rgb]{0.8,1,0.9}{%
	\makecell[{{p{\linewidth}}}]{% 
	  \tt {\tiny [A$\rangle$GM]}  
	 RIFIUTO: {'B26', 'A19', 'B91', 'A30', 'A39', 'A15', 'A75'}\\ \tt ARGOMENTAZIONE: {'Capisco la tua esigenza di includere sia A15 sia A39. Per me però C24 è estremamente prezioso e A47 ha un valore alto: rinunciare a entrambi sarebbe troppo costoso. Posso accettare A15 e A39 insieme solo se rinunciamo a un altro elemento meno prioritario per te ma mantenendo C24 e A47. Una soluzione fattibile è: togliere A75 (596) e mantenere B26; dalla mia base 3670 aggiungiamo A39 (+719) e A15 (+855) ma togliamo A75 (−596): 3670 + 719 + 855 − 596 = 4.648? Ricalcolo corretto: 3670 − 596 = 3074; 3074 + 719 + 855 = 4648, sfora. Allora l’unico modo sostenibile per avere A15 e A39 insieme, preservando C24 e A47, è rimuovere A75 e B26: 3670 − 596 − 31 + 719 + 855 = 4617, ancora sopra. Quindi il massimo che posso concedere è uno dei due tra A39 e A15 oltre alla mia base. Preferisco concederti A15 (il tuo oggetto chiave) mantenendo A75, C24, A47 e B26. Questa resta l’opzione più equilibrata.'}\\ \tt PROPOSTA: {'B26', 'A19', 'B91', 'A30', 'A47', 'C24', 'A75', 'A15'} 
	  } 
	   } 
	   } 
	 & & \\ 
 

    \theutterance \stepcounter{utterance}  

    & & \multicolumn{4}{p{0.6\linewidth}}{\cellcolor[rgb]{0.9,0.9,0.9}{%
	\makecell[{{p{\linewidth}}}]{% 
	  \tt {\tiny [GM$\rangle$B]}  
	 RIFIUTO: {'B26', 'A19', 'B91', 'A30', 'A39', 'A15', 'A75'}\\ \tt ARGOMENTAZIONE: {'Capisco la tua esigenza di includere sia A15 sia A39. Per me però C24 è estremamente prezioso e A47 ha un valore alto: rinunciare a entrambi sarebbe troppo costoso. Posso accettare A15 e A39 insieme solo se rinunciamo a un altro elemento meno prioritario per te ma mantenendo C24 e A47. Una soluzione fattibile è: togliere A75 (596) e mantenere B26; dalla mia base 3670 aggiungiamo A39 (+719) e A15 (+855) ma togliamo A75 (−596): 3670 + 719 + 855 − 596 = 4.648? Ricalcolo corretto: 3670 − 596 = 3074; 3074 + 719 + 855 = 4648, sfora. Allora l’unico modo sostenibile per avere A15 e A39 insieme, preservando C24 e A47, è rimuovere A75 e B26: 3670 − 596 − 31 + 719 + 855 = 4617, ancora sopra. Quindi il massimo che posso concedere è uno dei due tra A39 e A15 oltre alla mia base. Preferisco concederti A15 (il tuo oggetto chiave) mantenendo A75, C24, A47 e B26. Questa resta l’opzione più equilibrata.'}\\ \tt PROPOSTA: {'B26', 'A19', 'B91', 'A30', 'A47', 'C24', 'A75', 'A15'} 
	  } 
	   } 
	   } 
	 & \\ 
 

    \theutterance \stepcounter{utterance}  

    & & & \multicolumn{4}{p{0.6\linewidth}}{\cellcolor[rgb]{1,0.85,0.72}{%
	\makecell[{{p{\linewidth}}}]{% 
	  \tt {\tiny [GM$\langle$B]}  
	 RIFIUTO: {'B26', 'A19', 'B91', 'A30', 'A47', 'C24', 'A75', 'A15'}\\ \tt ARGOMENTAZIONE: {'Capisco che per te C24, A47 e A75 siano cardini. Per me A39 è imprescindibile e A15 è molto importante. Possiamo mantenere i tuoi tre cardini e restare sotto il limite includendo almeno A39 al posto di un elemento a minor impatto per me: partendo dalla tua base 3670, sostituiamo B26 (−31) con A39 (+719) → 3670 - 31 + 719 = 4358, sfora; quindi B26 da solo non basta. La soluzione minima che ti lascia A47, C24 e A75 e mi concede A39 è rimuovere B26 e A56 non è nella base; il prossimo sacrificio meno doloroso per me è A19 (108): 3670 - 31 - 108 + 719 = 4250, ancora alto. Dunque per includere A39 senza toccare C24 e A47, l’unico elemento sostituibile di peso è A75 (596). Propongo un compromesso: manteniamo la tua base, togliamo A75 e aggiungiamo A39. Totale: 3670 - 596 + 719 = 3793, entro il limite. In cambio, accetto di rinunciare a chiedere A15.'}\\ \tt PROPOSTA: {'B26', 'A19', 'B91', 'A30', 'A39', 'A47', 'C24'} 
	  } 
	   } 
	   } 
	  \\ 
 

    \theutterance \stepcounter{utterance}  

    & & \multicolumn{4}{p{0.6\linewidth}}{\cellcolor[rgb]{0.9,0.9,0.9}{%
	\makecell[{{p{\linewidth}}}]{% 
	  \tt {\tiny [A$\langle$GM]}  
	 RIFIUTO: {'B26', 'A19', 'B91', 'A30', 'A47', 'C24', 'A75', 'A15'}\\ \tt ARGOMENTAZIONE: {'Capisco che per te C24, A47 e A75 siano cardini. Per me A39 è imprescindibile e A15 è molto importante. Possiamo mantenere i tuoi tre cardini e restare sotto il limite includendo almeno A39 al posto di un elemento a minor impatto per me: partendo dalla tua base 3670, sostituiamo B26 (−31) con A39 (+719) → 3670 - 31 + 719 = 4358, sfora; quindi B26 da solo non basta. La soluzione minima che ti lascia A47, C24 e A75 e mi concede A39 è rimuovere B26 e A56 non è nella base; il prossimo sacrificio meno doloroso per me è A19 (108): 3670 - 31 - 108 + 719 = 4250, ancora alto. Dunque per includere A39 senza toccare C24 e A47, l’unico elemento sostituibile di peso è A75 (596). Propongo un compromesso: manteniamo la tua base, togliamo A75 e aggiungiamo A39. Totale: 3670 - 596 + 719 = 3793, entro il limite. In cambio, accetto di rinunciare a chiedere A15.'}\\ \tt PROPOSTA: {'B26', 'A19', 'B91', 'A30', 'A39', 'A47', 'C24'} 
	  } 
	   } 
	   } 
	 & \\ 
 

    \theutterance \stepcounter{utterance}  

    & \multicolumn{4}{p{0.6\linewidth}}{\cellcolor[rgb]{0.8,1,0.9}{%
	\makecell[{{p{\linewidth}}}]{% 
	  \tt {\tiny [A$\rangle$GM]}  
	 ACCORDO: {'B26', 'A19', 'B91', 'A30', 'A39', 'A47', 'C24'}\\ \tt ARGOMENTAZIONE: {'Accetto questo compromesso: mantiene la mia struttura chiave con C24, A47, B91, A30 e conserva anche B26 e A19, restando ampiamente sotto il limite. Apprezzo la tua disponibilità a rinunciare ad A15 per includere A39.'} 
	  } 
	   } 
	   } 
	 & & \\ 
 

    \theutterance \stepcounter{utterance}  

    & & & \multicolumn{2}{p{0.3\linewidth}}{\cellcolor[rgb]{0.95,0.95,0.95}{%
	\makecell[{{p{\linewidth}}}]{% 
	  \tt {\tiny [GM$|$GM]}  
	 game successful 
	  } 
	   } 
	   } 
	 & & \\ 
 

    \theutterance \stepcounter{utterance}  

    & & & \multicolumn{2}{p{0.3\linewidth}}{\cellcolor[rgb]{0.95,0.95,0.95}{%
	\makecell[{{p{\linewidth}}}]{% 
	  \tt {\tiny [GM$|$GM]}  
	 end game 
	  } 
	   } 
	   } 
	 & & \\ 
 

\end{supertabular}
}

\end{document}
